% ===== tr-05 块2c：第 4 章 pp.48--51（原文页 48--51）=====
% [原文 p.48]
% [Tutor Note: 编号提醒------「定理 4.10.1」来自原文 Theorem 10.1，「式 (4.10.1)」来自原文编号公式 (10.1)（p.51）；两者数字相同但分属定理与公式，这是原文按节编号的体系所致，并非重复编号，正文照此排。]
单纯性的定义相当繁复，故我们不在此给出；但它本质上说的是：不能有太多 $\check X(B)$（或 $X(B)$）的拓扑藏在 $B$ 的奇异轨迹之上。

[31] 的一个主要结果（定理 5.4）是

\begin{theorem}{4.10.1}
给定单纯的 $(B,\mathcal{P},\varphi)$，以 $(B,\mathcal{P},\varphi)$ 为相交复形的对数 Calabi--Yau 空间在「保持 $B$ 的同构」下的集合（即不交换不可约分支）是 $H^1(B,i_*\check\Lambda\otimes\mathbb{C}^\times)$。若一个同构在相交复形上诱导恒等，则称它保持 $B$。
\end{theorem}

于是模空间是一个代数环面（或代数环面的不交并），其维数等于 $\dim_{\mathbb{C}}H^1(B,i_*\check\Lambda\otimes\mathbb{C})$。在 [32] 中我们事实上证明了：这个环面的维数等于光滑纤维 $\mathcal{X}_t$ 的 $\dim H^1(\mathcal{X}_t,\mathcal{T}_{\mathcal{X}_t})\cong H^{n-1,1}(\mathcal{X}_t)$，其中 $\mathcal{X}_t$ 是 $X_0(B,\mathcal{P},s)^\dagger$ 的某个光滑化 $\mathcal{X}\to D$ 的光滑纤维。这正是预期的维数，因为后一个向量空间是 $\mathcal{X}_t$ 模空间的切空间。

现在假设给定一个对数 Calabi--Yau 空间 $X_0:=X_0(B,\mathcal{P},s)^\dagger\to 0^\dagger$。我们的目标是用对数结构提供「初始条件」，逐阶地产生 $k$ 阶形变 $X_k\to\operatorname{Spec}\mathbb{C}[t]/(t^{k+1})$。为此，我们将把 $X_0$ 的「各片段」的标准加厚粘合起来，并用一套我们称之为\emph{结构}（structure）的复杂数据来修改标准的粘合。

首先，我们所考虑的 $X_0$ 的「片段」是 $X_0$ 的各分层的环面开仿射子集。回忆 $X_0$ 的分层由胞腔 $\tau\in\mathcal{P}$ 标号，对应于一个射影环面簇 $\mathbb{P}_\tau$。还要回忆：若 $\omega\subseteq\tau$，则 $\tau$ 沿 $\omega$ 的法锥是定义 $\mathbb{P}_\tau$ 的扇中的一个锥，从而定义 $\mathbb{P}_\tau$ 的一个开仿射子集。我们把这个开仿射子集记作 $V_{\omega,\tau}\subseteq\mathbb{P}_\tau$；注意
\[
V_{\omega,\tau}=\mathbb{P}_\tau\setminus\bigcup_{\substack{\rho\subseteq\tau\\ \omega\not\subseteq\rho}}\mathbb{P}_\rho.
\]
例如，若 $\omega$ 是 $\tau$ 的顶点，则 $V_{\omega,\tau}$ 是 $\mathbb{P}_\tau$ 中包含对应于 $\omega$ 的零维分层的那个标准环面开仿射子集。

其次，集合 $V_{\omega,\tau}$ 的加厚是什么？它们可以如下显式地描述。选取 $\omega$ 内部一个不在 $B$ 的奇异轨迹 $\Gamma$ 中的点 $x$。我们在切空间 $\Lambda_x\otimes_{\mathbb{Z}}\mathbb{R}$ 中得到一个由未必严格凸的锥构成的扇 $\Sigma_x$，它由包含 $\omega$ 的每个胞腔 $\sigma$ 在 $x$ 处的切锥组成。我们可以选取 $\varphi$ 在 $x$ 的小邻域内的一个代表 $\varphi_x$，它在 $\omega$ 上为零；然后可以在 $\Sigma_x$ 的每个锥上线性延拓它，从而把 $\varphi_x$ 看作分段线性函数 $\varphi_x:\Lambda_x\otimes\mathbb{R}\to\mathbb{R}$。这又定义一个幺半群
\[
P_x:=\{(m,r)\in\Lambda_x\times\mathbb{Z}\mid r\ge\varphi_x(m)\},
\]
它与 $P_v$ 的定义完全类似。
% [原文 p.49]
对每个包含 $\tau$ 的极大胞腔 $\sigma$，令 $n_\sigma\in\check\Lambda_x$ 表示 $\varphi_x$ 限制到 $\sigma$ 的切锥上的斜率。于是我们在环 $\mathbb{C}[P_x]$ 中定义由下式给出的单项式理想
\[
I_{\omega,\tau}^{>k}=\big\langle z^{(m,r)}\ \big|\ (m,r)\in P_x,\ r-\langle n_\sigma,m\rangle>k\ \text{对某个}\ \sigma\in\mathcal{P}_{\max},\ \sigma\supseteq\tau\big\rangle.
\]
于是 $V_{\omega,\tau}$ 的所需标准加厚是
\[
V_{\omega,\tau}^k:=\operatorname{Spec}\mathbb{C}[P_x]/I_{\omega,\tau}^{>k}.
\]
容易验证：若 $k=0$，这就还原出 $V_{\omega,\tau}$；而若 $k>0$，则 $V_{\omega,\tau}^k$ 的约化空间是 $V_{\omega,\tau}$。因此这确实是 $V_{\omega,\tau}$ 的一个加厚。

有一点我们必须十分小心。这个定义表面上依赖于点 $x$，而由于奇异轨迹的存在，不同切空间 $\Lambda_x,\Lambda_{x'}$ 之间经由平行移动的等同依赖于路径。我们处理这一问题的办法不是选取特定的点 $x$，而是选取包含 $\tau$ 的一个特定的极大参考胞腔 $\sigma$。于是我们可以把任一 $\Lambda_x$ 与 $\Lambda_\sigma$（$\sigma$ 的、定义良好的切空间）经由从 $x$ 直接到 $\sigma$ 的平行移动等同起来。我们将把这一额外的参考胞腔选择记为 $V_{\omega,\tau,\sigma}^k$。不同的参考胞腔选择 $\sigma'$ 给出的空间 $V_{\omega,\tau,\sigma'}^k$ 与 $V_{\omega,\tau,\sigma}^k$ 抽象地同构，但不是典范同构。这一点在下面将会变得重要。我们还对坐标环使用记号
\[
R_{\omega,\tau,\sigma}^k:=\mathbb{C}[P_x]/I_{\omega,\tau}^{>k},
\]
同样要记住这一参考胞腔的选择。

这些各种加厚的概形之间也有自然的粘合。注意：给定 $\tau_1\subseteq\tau_2\subseteq\tau_3$，有自然的满射
\[
R_{\tau_1,\tau_3,\sigma}^k\to R_{\tau_1,\tau_2,\sigma}^k,
\]
给出闭嵌入 $V_{\tau_1,\tau_2,\sigma}^k\to V_{\tau_1,\tau_3,\sigma}^k$；以及自然的单射
\[
R_{\tau_1,\tau_3,\sigma}^k\to R_{\tau_2,\tau_3,\sigma}^k,
\]
给出开嵌入 $V_{\tau_2,\tau_3,\sigma}^k\to V_{\tau_1,\tau_3,\sigma}^k$。

若 $B$ 没有奇点，则参考胞腔 $\sigma$ 无关紧要，此时我们把它从记号中略去。特别地，容易验证：若取（比如说）$\tau_1$ 为固定的顶点 $v$，并取概形有向系 $\{V_{v,\tau}^k\mid v\in\tau\}$ 的极限，我们就得到 $V(v)$ 的一个 $k$ 阶加厚 $V^k(v)$，它由 $V^k(v)=U(v)\times_{\mathbb{A}^1}\operatorname{Spec}\mathbb{C}[t]/(t^{k+1})$ 给出，其中 $U(v)\to\mathbb{A}^1$ 是由 $z^{(0,1)}$ 给出的态射。这正是对数结构让我们期待的那种最朴素的光滑化。注意我们可以把这个概形的正向极限写成
\[
\operatorname{Spec}\varprojlim_{\tau}R_{v,\tau}^k.
\]

基本想法于是将是：用某些额外数据来修改上面的各种映射。
% [原文 p.50]
\figimg[0.72]{p4_fig05.png}
\figtitle{4.5}{基本例子。该图显示两个坐标卡的仿射嵌入，它们由按两种不同方式切割两个三角形之并而得到。每个三角形都是一个标准单纯形。}

为了理解为什么需要这些修改，让我们考虑最重要的单个例子，即二维 $B$ 中一个 focus-focus 型的孤立奇点。

我们假设 $\mathcal{P}$ 含两个极大胞腔 $\sigma_1,\sigma_2$，且 $\sigma_1\cap\sigma_2=\tau$，如图 4.5 所示。注意两个坐标卡的交是 $(\sigma_1\cup\sigma_2)\setminus\tau$，而过渡映射在 $\sigma_1\setminus\tau$ 上是恒等，在 $\sigma_2\setminus\tau$ 上由线性变换
% [Tutor Note: 这里过渡映射的矩阵是 \begin{pmatrix}1&0\\1&1\end{pmatrix}，与例 4.7.6 中绕奇点的单值矩阵 \begin{pmatrix}1&4\\0&1\end{pmatrix} 不同类；前者是两坐标卡之间的粘合变换，后者是单值变换，切勿互换。]
\[
\begin{pmatrix} 1 & 0 \\ 1 & 1 \end{pmatrix}
\]
给出。这两个坐标卡合起来在 $(\sigma_1\cup\sigma_2)\setminus\Gamma$ 上定义一个整仿射结构，其中 $\Gamma=\{p\}$ 是两处切割的公共点。

然后我们取 $\varphi$ 为单值的：在 $\sigma_1$ 上恒为 $0$，在右方顶点处取值 $1$。

于是可以算出
\[
\begin{aligned}
R_{\tau,\sigma_1,\sigma_1}^k&=\mathbb{C}[x,y,w^{\pm1}]/(y^{k+1})\\
R_{\tau,\sigma_2,\sigma_2}^k&=\mathbb{C}[x,y,w^{\pm1}]/(x^{k+1})\\
R_{\tau,\tau,\sigma_i}^k&=\mathbb{C}[x,y,w^{\pm1}]/(x^{k+1},y^{k+1}).
\end{aligned}
\]
这里，若我们使用左方的坐标卡，即选取 $p$ 下方的一点 $s$ 并在 $P_s\subseteq\Lambda_s\oplus\mathbb{Z}$ 中工作，则变量 $x,y$ 与 $w$ 与 $\mathbb{C}[P_s]$ 的元素如下等同：
\[
x=z^{(-1,0,0)},\qquad y=z^{(1,0,1)},\qquad w=z^{(0,1,0)}.
\]
我们有自然的满射 $R_{\tau,\sigma_i,\sigma_i}^k\to R_{\tau,\tau,\sigma_i}^k$，并通过经由 $s$ 的平行移动把 $\Lambda_{\sigma_1}$ 与 $\Lambda_{\sigma_2}$ 等同，从而把 $R_{\tau,\tau,\sigma_1}^k$ 与 $R_{\tau,\tau,\sigma_2}^k$ 等同起来。由于我们已把一切写在左方坐标卡中，而 $\Lambda_{\sigma_1}$ 与 $\Lambda_{\sigma_2}$ 在那里经由 $s$ 的平行移动被等同，这个等同就是平凡等同。于是我们可以把各加厚的坐标环粘合为
\[
R_{\tau,\sigma_1,\sigma_1}^k\times_{R_{\tau,\tau,\sigma_i}^k}R_{\tau,\sigma_2,\sigma_2}^k.
\]
这个环的纤维积容易看出同构于环
\[
\mathbb{C}[X,Y,W^{\pm1},t]/(t-XY,t^{k+1}),
\]
% [原文 p.51]
\figimg[0.72]{p4_fig06.png}
% [Tutor Note: 原文此图只有编号「FIGURE 6.」而无图注文字，故 \figtitle 第二参数留空；图内标注（上方射线 $1+w^{-1}$、下方射线 $1+w$）已含在裁图 p4_fig06.png 中。这两条射线的意义见正文：上方射线对应 p.52 的自同构 (10.3)（含 $1+w^{-1}$），下方射线对应 (10.2)（含 $1+w$）------(10.2)、(10.3) 属下一块 tr-06，编号应为 4.10.2、4.10.3。]
\figtitle{4.6}{}
其中 $X=(x,x)$、$Y=(y,y)$、$W=(w,w)$，而 $t=(xy,xy)$，都是作为环的笛卡尔积中的元素。

另一方面，假设我们改为经由位于 $p$ 上方的一点 $s'$ 的平行移动来等同 $R_{\tau,\tau,\sigma_1}^k$ 与 $R_{\tau,\tau,\sigma_2}^k$。为此我们可以在右方坐标卡中工作。同样，$x,y$ 与 $w$ 是用 $\sigma_1$ 中的切向量 $(-1,0),(1,0)$ 与 $(0,1)$ 定义的，而在第二个坐标卡中它们被移到 $\sigma_2$ 中的相同切向量。然而，为了把它与我们原先对 $R_{\tau,\tau,\sigma_2}^k$ 的描述相比较，我们需要把这些切向量看作 $\sigma_2$ 在原来坐标卡（即左方坐标卡）中的切向量。在那里，这些切向量分别是 $(-1,1),(1,-1)$ 与 $(0,1)$。于是我们得到一个同构 $R_{\tau,\tau,\sigma_1}^k\to R_{\tau,\tau,\sigma_2}^k$，它由
\begin{equation}
x\mapsto xw,\qquad y\mapsto yw^{-1},\qquad w\mapsto w \tag{4.10.1}
\end{equation}
给出。

利用这一等同，我们得到一个复合映射 $R_{\tau,\sigma_1,\sigma_1}^k\to R_{\tau,\tau,\sigma_1}^k\to R_{\tau,\tau,\sigma_2}^k$，从而得到纤维积
\[
R_{\tau,\sigma_1,\sigma_1}^k\times_{R_{\tau,\tau,\sigma_2}^k}R_{\tau,\sigma_2,\sigma_2}^k\cong\mathbb{C}[X,Y,W^{\pm1},t]/(XY-tW,t^{k+1}),
\]
其中现在
\[
X=(x,xw),\qquad Y=(yw,y),\qquad W=(w,w),\qquad t=(xy,xy).
\]
注意这个新环虽然抽象地同构于前一个环，但作为 $\mathbb{C}[t]/(t^{k+1})$-代数并不同构。

于是粘合定义得不妥，而这正是由 $B$ 的奇点造成的。此时正确的光滑化将依赖于对数结构的选择，但无论如何我们预期它应当是一个形如 $\operatorname{Spec}\mathbb{C}[X,Y,W^{\pm1},t]/(XY-f(W)t)$ 的族，其中 $f(W)$ 是某个沿 $W$-轴恰好在该给定 $X_0(B,\mathcal{P},s)$ 上对数结构不精细的那些点处消失的函数。显然 $f$ 于是被对数结构在相差可逆函数的意义下确定。为这个例子起见，让我们取函数 $f(W)=1+W$，注意 $f(W)=1+W^{-1}$ 也同样可以。现在我们可以借助图 4.6 来修改这些粘合。

在这幅图中，我们画出了包含在 $\tau$ 中、从奇点发出的两条射线，并分别把这两个箭头标以函数 $1+w^{-1}$ 与 $1+w$。这些射线告诉我们：若我们试图用两个极大胞腔之间的平行移动把 $R_{\tau,\tau,\sigma_1}^k$ 与 $R_{\tau,\tau,\sigma_2}^k$ 等同起来，就需要用一个由穿过其中一条射线所给出的自同构来修改这一等同。
